\section*{Reproducibility}

Every number in this paper comes from measurement files committed alongside
the code that produced them, and the manuscript is built from those files
rather than from transcribed values. The repository will be released under
the terms confirmed with the venue of the earlier version; see the footnote
in Section~\ref{sec:intro}.

The four experiments are run as follows, each through the wrapper that
enforces the pinned environment and exits non-zero if it is not present:

\begin{verbatim}
  bash tools/run.sh python -m experiments.e1_resolution_sweep
  bash tools/run.sh python -m experiments.e2_erf
  bash tools/run.sh python -m experiments.e3_coverage
  bash tools/run_e4.sh
\end{verbatim}

The first three are minutes to hours. The fourth is not: it launches twelve
training runs of three hundred epochs each and takes \HoursTotal{} GPU-hours,
about four days on the single device described in
Section~\ref{sec:protocol}. It writes a checkpoint every epoch and the same
command resumes an interrupted campaign, which we exercised twice. We state
the cost because a reproducibility claim that omits it is not one.

The tables, macros, figures, and PDF of this paper are regenerated from the
committed measurement files by:

\begin{verbatim}
  bash tools/run.sh python -m paper.make_tables
  bash tools/run.sh python -m paper.make_figures
  bash tools/latex.sh
\end{verbatim}

The test suite fails if a regenerated table differs from the committed one,
if the four experiments disagree about the environment they ran in, or if a
number that looks like a measurement appears in the prose without coming
from a macro.
